\documentclass[11pt]{article}
\usepackage[T1]{fontenc}
\usepackage{lmodern,microtype,amsmath,amssymb,amsthm,booktabs,array,graphicx,xcolor,hyperref,fancyhdr}
\usepackage[letterpaper,top=.85in,bottom=.84in,left=.94in,right=.94in,headheight=14pt]{geometry}
\definecolor{ink}{HTML}{182A35}\definecolor{teal}{HTML}{176662}
\hypersetup{colorlinks=true,linkcolor=teal,urlcolor=teal,citecolor=teal,pdftitle={Arithmetic criteria for three Riemann hypotheses},pdfauthor={Research report}}
\pagestyle{fancy}\fancyhf{}\lhead{\small Arithmetic criteria for three Riemann hypotheses}\rhead{\small Constructive arithmetic}\cfoot{\small\thepage}
\renewcommand{\headrulewidth}{.3pt}
\newtheorem{theorem}{Theorem}[section]\newtheorem{lemma}[theorem]{Lemma}
\newcommand{\code}[1]{\texttt{\detokenize{#1}}}
\newcommand{\N}{\mathbb N}\newcommand{\Q}{\mathbb Q}\newcommand{\Z}{\mathbb Z}
\newcommand{\pos}[1]{\left(#1\right)_{+}}
\newcommand{\sgap}{D_F}
\raggedbottom\urlstyle{same}
\setcounter{tocdepth}{1}
\begin{document}
\thispagestyle{plain}
{\small\sffamily\color{teal} ANALYTIC NUMBER THEORY \enspace/\enspace MATHEMATICAL LOGIC}\par\vspace{12pt}
{\LARGE\bfseries\color{ink} Arithmetic criteria for three Riemann hypotheses}\par\vspace{5pt}
{\large Sparse prime counts, Taylor coefficients, and rational certificates in iPRA}\par\vspace{6pt}
{\small November 20, 2026}
\bigskip

One finite computation at each of a sparse sequence of cutoffs suffices to express the Riemann hypothesis for \emph{every} number field. An explicit schedule selects one monic integer polynomial at each cutoff. It measures only an upward discrepancy in an integer built from rational prime powers. Each individual polynomial returns often enough to detect a zero anywhere off the critical line. Changing which polynomial is selected gives ordinary RH or Dirichlet GRH with the same algorithm.

\begin{center}\small
\begin{tabular}{@{}p{.15\linewidth}p{.48\linewidth}p{.28\linewidth}@{}}\toprule
Conjecture & Functions whose zeros in $0<\Re s<1$ must lie on $\Re s=1/2$ & Polynomials tested\\\midrule
RH & Riemann zeta function $\zeta(s)$ & $X-1$\\[3pt]
Dirichlet GRH & $L(s,\chi)$ for every primitive Dirichlet character $\chi$, including complex characters & $X^d-1$, $d\ge1$\\[3pt]
Dedekind ERH & $\zeta_K(s)$ for every number field $K$, including nonnormal fields & Every monic $F\in\Z[X]$, $\deg F\ge1$\\\bottomrule
\end{tabular}
\end{center}

A single Taylor expansion after a rational change of variable reduces all three to signed integer coefficient tests. Polynomial multiplication on Taylor disks leads to finite matrix inequalities. For ordinary RH, the Robin and Nicolas inequalities admit distinct weak and strict arithmetic readings, and the discrete Nyman--Beurling criterion admits rational least-squares witnesses with a certified tail.

\paragraph{What is proved where.}
The arithmetic statements are literal closed formulas in ProveIt's iPRA. Empty-context derivations in that calculus prove such equivalences as the two-sided prime test versus the conjunction of its one-sided tests, and strict divisor or totient certificates versus a least successful precision. Lean also checks the mathematical readings of finite computations. The identifications with zeros use classical analytic number theory, with their proofs and dependencies given below. The report marks which links have each strength. The conjectures themselves remain open.

\tableofcontents
\newpage
\section{The object language and its three number-theoretic sentences}\label{sec:foundation}
The ordinary first-order presentation of intuitionistic primitive recursive arithmetic (iPRA) has one numerical sort, a symbol for each particular primitive recursive program, equality, intuitionistic connectives, and number quantifiers. Induction is restricted to quantifier-free formulas, which may include defined primitive recursive symbols. The repository presents the syntax as \code{PR k}, \code{PTm n} and \code{Fm PTm n}. A \code{PR k} is built from zero, successor, projections, composition and primitive recursion. Free variables of a term or formula lie among $n$ de Bruijn indices. The type \code{Prf PRAAx [] A} expresses that the closed sentence $A$ is derivable without hypotheses by the repository's rules.

Real numbers, functions quantified as objects, and sets of zeros do not occur in this language. We name the three hypotheses by arithmetic sentences, prove transformations inside iPRA where available, and identify their meaning in the standard model. A classical truth equivalence alone gives no term of the type \code{Prf PRAAx [] (Fm.iff A B)}. In particular, a classical argument producing a positive margin or a precision by existence has an additional obligation if one asks for a primitive recursive witness transformation.

\paragraph{Family parameters.}
Let $d_0,H,u,v$ be naturals. Write $d=d_0+1$, $b=H+2$ and $[w]_i=\lfloor w/b^i\rfloor\bmod b$. They name
\[
 F(X)=X^d+\sum_{i<d}([u]_i-[v]_i)X^i.
\]
Every monic nonconstant polynomial occurs by choosing a base above all coefficient magnitudes. A change of base can change the allowance below. A fixed code is retained when a bound depends on that allowance.

At each prime power put
\[
 h_F(p,k)=\deg\gcd(\bar F,X^{p^k}-X)\quad \text{over }\mathbb F_p,\qquad
 B_F(y)=\left\lfloor\log_2 \prod_{p^k\le y}p^{h_F(p,k)}\right\rfloor .
\]
Here $k\ge1$. The product of no factors is one. The gcd counts each distinct irreducible factor of $\bar F$ of degree dividing $k$, weighted by its degree. This follows by factoring $X^{p^k}-X$ over $\mathbb F_p$ and observing that its derivative is $-1$. The PR code checks divisibility and irreducibility by bounded searches through coefficient words, including when $\bar F$ has repeated roots. The position of the highest set bit, numbered from zero, implements the binary logarithm rounded down. Put
\[
 L(y)=1+\lfloor\log_2 y\rfloor\quad(y\ge1),\quad C=1000(d+1)^4(H+2)^2,\quad
 D_F(n)=B_F(2n)-2B_F(n).
\]
The original matrix is the excess
\[
 E(d_0,H,u,v,t)=\max\{0,D_F(n)^2-nC^2L(2n)^4\},\qquad n=100+t.
\]
RH is $\forall t\ E(0,0,0,1,t)=0$. GRH also quantifies $d_0$, with $(H,u,v)=(0,0,1)$. ERH quantifies all five arguments. With $F=X-1$, $B_F(y)$ is the index of the highest set bit of $\mathrm{lcm}(1,\ldots,\lfloor y\rfloor)$. These $\Pi^0_1$ formulas admit decidable finite failures. iPRA proves $\mathrm{ERH}\Rightarrow\mathrm{GRH}\Rightarrow\mathrm{RH}$ by universal specialization. For each it also proves equivalence to the negation of an existential failing tuple. In contrast, passing from the negation of a universal test to the existence of such a tuple would require an additional principle in general.

\paragraph{Why these integers recover the required functions.}
Factor $F=\prod_i f_i^{m_i}$ over $\Q$ with $f_i$ distinct and set $K_i=\Q[X]/f_i$. Let $G=\prod_i f_i$, $\Delta=|\mathrm{disc}(G)|$ and
$\Psi_F(y)=\sum_{p^k\le y}h_F(p,k)\log p$. Write $r$ for the number of factors. If $p\nmid\Delta$, reduction of the order generated by a root of $f_i$ gives the same finite algebra as reduction of its ring of integers. Distinct factors stay coprime and unramified. The residue degrees above $p$ therefore give exactly $h_F(p,k)$ when summed over degrees dividing $k$. At the exceptional primes either count is between $0$ and $d$. Consequently
\[
 \left|\Psi_F(y)-\sum_i\psi_{K_i}(y)\right|
 \le d\,\omega(\Delta)\log y,\qquad
 \sum_i\log |\Delta_{K_i}|\le\log\Delta\le 2d^2(H+2).
\]
The last inequalities follow from the discriminant of a power basis, the product formula for discriminants and resultants, and the root bound $|\alpha|\le H+2$. They hold in degree one with $\Delta=1$.

Under the zero hypothesis for the $K_i$, the explicit prime ideal theorem of Greni\'e and Molteni (Corollary 1.2) bounds $|\psi_K(y)-y|$ for $y\ge100$ by
\[
 \sqrt y\left[\left(\frac{\log y}{2\pi}+2\right)\log|\Delta_K|
          +\left(\frac{\log^2 y}{8\pi}+2\right)[K:\Q]\right].
\]
Adding over $i$ and using the exceptional-prime bound gives
$|\Psi_F(y)-ry|\le13d^3(H+2)\sqrt y\,L(y)^2$ at integer $y\ge100$.
Indeed each parenthesis is at most $3L(y)^2$ and $\omega(\Delta)\le2\log\Delta$.
The linear terms cancel at $2n$ and $n$. Flooring after division by $\log 2$ costs at most three bits, which is absorbed in $C\sqrt n L(2n)^2$. Thus every arithmetic test holds under the corresponding hypothesis. The converse follows already from the upper half alone by Theorem~\ref{thm:oneside}.

\paragraph{From the fields to the named hypotheses.}
Every number field has an integral primitive element, so its minimal polynomial is an allowed input. For $F=X^d-1$ the distinct factors are $\Phi_e$, $e\mid d$. The field $\Q(\zeta_e)$ has zeta function equal to the product of primitive Dirichlet $L$-functions of conductor dividing $e$, including conductor one. In the open strip all factors are holomorphic, so their zeros persist in the product with multiplicity. Taking $d$ equal to the conductor reaches each primitive character. The functional equations reflect nontrivial zeros across the critical line, for conjugate characters when necessary. Boundary zeros of imprimitive Euler factors on $\Re s=0$ have no role. For the terminology see the articles on the \href{https://en.wikipedia.org/wiki/Riemann_hypothesis}{Riemann hypothesis} and the \href{https://en.wikipedia.org/wiki/Generalized_Riemann_hypothesis}{generalized Riemann hypothesis}. For the analytic inputs see the \href{https://arxiv.org/abs/1312.4463}{explicit prime ideal theorem} and \href{https://math.mit.edu/classes/18.785/2018fa/LectureNotes19.pdf}{Sutherland, Theorem 19.15}.

\section{An upper prime discrepancy on sparse cutoffs}\label{sec:sparse}
For a fixed polynomial let $D_F(n)=B_F(2n)-2B_F(n)$ as in Section~\ref{sec:foundation}. The two-sided test treats upward and downward fluctuations alike. The following result identifies a sign that works for number fields even in the presence of a possible real zero. Write $x_+=\max(0,x)$.

\begin{theorem}\label{thm:oneside}
Fix a code for $F$. The critical-strip zeros of every field factor of $F$ are on the critical line if and only if
\[
     D_F(n)_+^2\le nC^2L(2n)^4\qquad(n\ge100).
\]
The reverse implication holds with \emph{any} fixed finite coefficient in place of $C$ and with any finite initial segment omitted. For $F=X-1$ the same conclusions hold with $(-D_F(n))_+$.
\end{theorem}
\begin{proof}
The forward implication follows from the two-sided estimate in Section~\ref{sec:foundation}. For the converse write $V(x)=\Psi_F(2x)-2\Psi_F(x)$ for real $x\ge1$. Passing from $D_F(n)$ to $V(n)$ changes the quantity by a factor $\log 2$ and at most $3\log 2$ from the two floors. On the interval $n\le x<n+1$, $\Psi_F(x)$ is constant and $\Psi_F(2x)$ has at most one extra jump, of size at most $d\log(2n+1)$. A jump can occur only at a prime power. Hence for some constant $K>0$ depending on the fixed data we have
\[
     V(x)\le W(x):=K\sqrt{x}(1+\log x)^2\qquad(x\ge1).
\]
Increasing $K$ handles any initial omissions. Define $J_F(s)=\sum_{p,k} h_F(p,k)\log p\, p^{-ks}$. For $\Re s>1$, integration of the step functions gives
\begin{align*}
 \int_1^\infty V(x)x^{-s}\,\frac{dx}{x}&=\frac{2^s-2}{s}J_F(s),\\
 J_F(s)&=-\sum_i\frac{\zeta_{K_i}'(s)}{\zeta_{K_i}(s)}+U_F(s).
\end{align*}
Absolute convergence justifies both the change of scale and termwise integration, since the coefficient at an integer $n$ is at most $d\log n$. The lower endpoint causes no term because $\Psi_F(x)=0$ for $1\le x<2$. The good primes give the displayed logarithmic derivatives. The correction $U_F$ is holomorphic on $\Re s>0$ because its coefficients at each of finitely many primes are bounded in modulus by $d\log p$ and the resulting power series converge geometrically there.

Take the Mellin transform $A(s)$ of the nonnegative function $W-V$. Its transform converges absolutely for $\Re s>1$. It admits the meromorphic continuation
\[
 A(s)=K\left(\frac1z+\frac2{z^2}+\frac2{z^3}\right)
      -\frac{2^s-2}{s}J_F(s),\qquad z=s-\tfrac12.
\]
There is no pole at $1$ since the zero of $2^s-2$ cancels the simple logarithmic-derivative poles. The zeta factors have no zeros on the real ray $s>1$ by their Euler products. Any real pole at $\beta\in(1/2,1)$ must come from zeros of total multiplicity $m_\beta>0$. Its residue in $A$ is
\[
                   \frac{(2^\beta-2)m_\beta}{\beta}<0.
\]
The sign matters here. Each field factor contributes the same sign, including when the same field occurs twice.

Let $\sigma_c$ be the real abscissa of convergence for the nonnegative integral. If it is $-\infty$ the integral is entire. If it is finite, Landau's lemma makes $\sigma_c$ a singularity. To see the lemma directly, put $x=e^t$. If the integral had an analytic continuation near $\sigma_c$, take a real center $a$ just to its right so that its open Taylor disk reaches a real point $a-h<\sigma_c$. The disk exists by choosing $a$ sufficiently near the center of the hypothetical analytic neighborhood. For the expansion in $a-s$, every coefficient is the integral of $(W(e^t)-V(e^t))e^{-at}t^k/k!$, which is nonnegative. Differentiation is justified using the convergent integral at a real number between $a$ and $\sigma_c$. Summing at $h$ and using monotone convergence gives the finite integral at $a-h$, a contradiction.

If $\sigma_c>1/2$, the meromorphic formula and $\sigma_c\le1$ now force a simple real pole at $\sigma_c$ with negative residue. Along real $s\downarrow\sigma_c$ this makes $A(s)$ tend to $-\infty$, impossible for an integral of a nonnegative function. Therefore $\sigma_c\le1/2$. Absolute convergence locally in $\Re s>1/2$ makes $A$ holomorphic there. A zero $\rho$ with $1/2<\Re\rho<1$ would leave a pole, because $|2^\rho|<2$ and the logarithmic residues add. The functional equations give the other half of the strip.

For the lower bound on ordinary $\zeta$, use the nonnegative function $W+V$. There are no real zeta zeros in $(1/2,1)$: the alternating eta series is strictly positive there by pairing consecutive terms and $1-2^{1-s}$ is strictly negative. Thus the real abscissa argument applies even with the sign reversed. Extending this last proof to an arbitrary field would need a way of handling its possible real zeros. The upper-bound proof already handles those zeros.
\end{proof}

\subsection*{Testing a single polynomial at a given cutoff}
Let $\nu_p(m)$ be the exponent of $p$ in a positive integer $m$. It is a PR function, given by a maximum over $0\le e\le m$ with $p^e\mid m$. For $j\ge0$ define
\[
 Q_j=\sum_{1\le i\le j} L(i),\qquad n_j=(10+Q_j)^2.
\]
Thus $Q_j$ counts the bits in the ordinary binary representations of $1,\ldots,j$, and $n_0=100$. Choose the polynomial data at stage $j$ by the following table. In each entry the valuations are taken at $j+1$.
\begin{center}
\begin{tabular}{@{}lcccc@{}}\toprule
Family &$d_0$&$H$&$u$&$v$\\\midrule
RH&0&0&0&1\\
Dirichlet GRH&$\nu_2$&0&0&1\\
Dedekind ERH&$\nu_2$&$\nu_3$&$\nu_5$&$\nu_7$\\\bottomrule
\end{tabular}
\end{center}
The entry for ERH with all four coordinates zero codes $F=X$; the ordinary RH specialization uses $F=X-1$. Denote the selected polynomial by $F_j$ and its particular allowance by $C_j$.
\begin{theorem}\label{thm:rotate}
For each row the corresponding analytic conjecture is equivalent in the standard model to the single closed $\Pi^0_1$ sentence
\begin{equation}\label{eq:rotating}
 \boxed{\quad\forall j\in\N\qquad D_{F_j}(n_j)_+^2\le n_j\,C_j^2 L(2n_j)^4.\quad}
\end{equation}
The same assertion with $n_j=(10+j)^2$ or $n_j=100+j$ is equivalent as well. Each version visits just one polynomial at a given stage. For ordinary RH either sign suffices on each of the three grids.
\end{theorem}
\begin{proof}
Under the relevant zero hypothesis each row satisfies the original bound at every cutoff. It remains to prove that visiting many fewer cutoffs and field codes still detects every offending field.

Fix an ERH code $(a,b,c,e)$, put $M=2^a3^b5^c7^e$ and $P=210M$. Each stage $j_h=M(1+210h)-1$ has exactly the required valuations. Indeed $1+210h$ is coprime to $2,3,5,7$. Consecutive such stages are exactly $P$ indices apart. For the Dirichlet row one can simply take the subsequence with the other three exponents zero. For ordinary RH all stages test the required polynomial. No scheduling constant uniform in the field is needed for the converse.

Set $r_j=10+Q_j$. On the displayed progression,
\[
 r_{j+P}-r_j\le P L(j+P),\qquad
 n_{j+P}-n_j\le 2r_jP L(j+P)+P^2L(j+P)^2 .
\]
For fixed $P$ and large $j$, $r_j\ge j$, $L(j+P)=O_P(1+\log r_j)$, and the right side is $O_P(\sqrt{n_j}(1+\log n_j))$. In particular $n_{j+P}/n_j\to1$. The same mesh bound holds for the ordinary square grid and the linear grid by taking their differences explicitly.

Take any real $x$ between two successive cutoffs $N,N'$ for this fixed code. The integer jumps of $\Psi_F$ are bounded by $d\log(2N')$ up to $2N'$, whence
\[
 |V(x)-V(N)|\le 4d(N'-N+1)\log(2N').
\]
This follows by counting the possible jumps on $(2N,2x]$ and twice those on $(N,x]$. The mesh estimate makes this error $O_F(\sqrt{x}(1+\log x)^2)$. Thus the upper bound at the selected nodes yields the upper envelope $V(x)\le K_F\sqrt x(1+\log x)^2$ for all large $x$. The proof of Theorem~\ref{thm:oneside} applies. The same interpolation is valid for the ordinary lower test. Since $F$ was arbitrary, Section~\ref{sec:foundation} gives precisely the indicated family.
\end{proof}

The number of sampled cutoffs below $X$ is asymptotic to $2\sqrt X/\log_2 X$, since $Q_j=j\log_2 j+O(j)$ by summing $\log_2 i+O(1)$. This counts evaluations only. The PR program at a particular cutoff can be very expensive. A fixed offending field would force arbitrarily large positive values of $D_F(n_j)/(\sqrt{n_j}L(2n_j)^2)$ even along its prescribed return progression: otherwise the same interpolation and transform proof would exclude its zero. This argument proves existence in classical analysis and supplies no bound on the first violating index.

\paragraph{Exact check of the schedules.}
The first five $n_j$ are $100,121,169,225,324$. For $F=X-1$ the corresponding signed discrepancies are $27,8,-7,28,14$. At the third cutoff the upper test sees zero positive discrepancy. At that same index the ERH scheduler selects $(0,1,0,0)$ and hence $F=X$. The script \code{check_sampling.py} reproduces the values from lcms and checks specified returns for 256 distinct four-word codes.

\paragraph{Object-theoretic content.}
\code{OneSided.lean} uses truncated subtraction to compute $p=(B_F(2n)-2B_F(n))_+$ and routes the original excess to one of two cells. For $p=0$ these cells are $(0,E)$, and for $p>0$ they are $(E,0)$. At a tie $E=0$. Induction on $p$ with a quantifier-free invariant proves that $E=0$ is equivalent to both cells being zero. Lifting across all family binders gives the object theorem \code{OneSided.split}. Specializing those binders at the scheduled valuation terms proves \code{to_rotating}. The file also constructs derivations connecting finite failures to the one-parameter tests and checks their arithmetic interpretation in the natural numbers. The accompanying \code{GridInterpolation.lean} identifies the compiled valuation with prime factorization, verifies the four-prime return progression and the mesh estimate, and proves a signed interpolation inequality from bounds on individual jumps. The declarations \code{compiled_log_grid} and \code{selectedData_return} connect the compiled sample to that grid and the return progression. The converse from the sparse upper test to the full two-sided test uses the preceding Landau argument. Its proof here is classical.

\section{One Taylor center and one sign suffice}\label{sec:cayley}
An ordinary Taylor disk centered at $5$ misses points just to the right of
the critical line at large height. A rational change of variable sends
\emph{the entire right half-plane} to the unit disk. At every positive
real radius the transformed function can already be evaluated by an
absolutely convergent series. This extra information lets us discard one
sign of every coefficient. We give closed integer tests, including a version
requiring only one precision per coefficient.

\subsection*{The data and the six tests per family}
Let $F$ be the coded monic polynomial of degree $d\ge1$, with the height
parameter $H$ of the input packet, and put $C=1000(d+1)^4(H+2)^2$.
For reference, write
\[
 h_F(p,k)=\deg\gcd(\bar F,X^{p^k}-X),\quad
 B(y)=\left\lfloor\log_2\prod_{p^k\le y}p^{h_F(p,k)}\right\rfloor,
 \quad a_n=B(n)-B(n-1).
\]
Primes $p$ and positive integers $k$ occur in the product, with $B(0)=0$.
Define $b_n=a_n-2a_{n/2}$ at even $n$, and $b_n=a_n$ at odd $n$.
For a cutoff $N\ge1$ let
\[
 J_n=\sum_{i=1}^{N(n-1)}\left\lfloor\frac{N^2}{N+i}\right\rfloor,
 \qquad A_k=\sum_{n=1}^N b_n J_n^k (N!/n)^5,
 \qquad D_k=k!N^{2k}(N!)^5 .
\]
All denominators below are positive. At each order $m\ge1$ put
\begin{align*}
 T_{m,N}
   &=\sum_{k=1}^m (-9)^k\binom{m-1}{k-1}
       \frac{m!}{k!}N^{2(m-k)} A_k,\\
 E_m&=81(d+1)10^{m-1},\qquad V_m=13500C(m+1)^4 .
\end{align*}
Choose \emph{one} of the upper, lower and two-sided comparisons
\begin{equation}\label{eq:cayleytest}
  NT_{m,N}\le D_m(N V_m+E_m),\qquad
  -NT_{m,N}\le D_m(N V_m+E_m),\qquad
  N|T_{m,N}|\le D_m(N V_m+E_m).
\end{equation}
Quantify over every $m\ge1$. Independently choose between testing
\emph{every} $N\ge1$ and testing only $N=E_m$.
For RH use $F=X-1$, for GRH use every $F=X^q-1$ with
$q\ge1$, and for ERH use every polynomial packet. Thus no zero, contours,
real constants, choice functions, or numerical witnesses occur in a test.
A signed inequality is compiled by moving its two nonnegative sums to
opposite sides and testing whether their truncated difference is zero.

\paragraph{Theorem.}
Each of these six universal sentences, interpreted in the standard
integers, is equivalent to the corresponding analytic hypothesis. A single
sign at the prescribed precision is sufficient. In the two-sided direction
the coefficient estimate loses at most one power of $m+1$ when recovered
from one sign. The explicit bound is given below.

\subsection*{From off-line zeros to poles of one germ}
Use $H_F(s)=\sum_{n\ge1} b_n n^{-s}$ for $\Re s>1$.
We recall the analytic inputs from the attached prime-product construction.
If $J_F(s)=\sum_{p,k}h_F(p,k)(\log p)p^{-ks}$ and
$\Psi_F(x)=\sum_{p^k\le x}h_F(p,k)\log p$, Abel summation gives
\[
 H_F(s)=(1-2^{1-s})\left(\frac{J_F(s)}{\log 2}
 +s\int_1^\infty
 \left(B(\lfloor x\rfloor)-\frac{\Psi_F(x)}{\log 2}\right)x^{-s-1}\,dx\right).
\]
The discrepancy inside the integral is bounded by $1$, so the integral
is holomorphic on $\Re s>0$. Outside finitely many bad primes, the local
factors identify $J_F$ with minus the sum of logarithmic derivatives of
the Dedekind zeta functions belonging to the distinct irreducible factors
of $F$ over $\mathbb Q$. At a bad prime the correction is holomorphic on
$\Re s>0$ by its geometric-series expression. At a common zero the
residues of the logarithmic derivatives add with the same sign. In
$1/2<\Re s<1$, the factor $1-2^{1-s}$ is nonzero; it cancels the
possible pole at $s=1$. Consequently an off-line zero on the right gives
a genuine pole. The functional equations reflect zeros on the left.
The cyclotomic factorization accounts for all primitive Dirichlet
characters, and an integral primitive element gives a packet for every
number field. These are classical identification steps.

The other input is the numerical bound under the corresponding hypothesis,
\begin{equation}\label{eq:cayleysum}
   \left|\sum_{n\le x} b_n\right|\le 100 C\sqrt{x}(1+\log x)^2
   \qquad(x\ge1).
\end{equation}
Indeed the prefix is $B(\lfloor x\rfloor)-2B(\lfloor x/2\rfloor)$.
At an even cutoff $2n\ge200$ this is the original dyadic gap with bound
$C\sqrt n(1+\lfloor\log_2(2n)\rfloor)^2$. An adjacent odd cutoff
adds at most $d\log_2(2n+1)+1$. The elementary bound
$B(y)\le d y\log_2 y$ covers $1\le y<200$. These give
\eqref{eq:cayleysum} with the displayed constant.

Use the change of coordinate
\[
 s(z)=\frac{5+4z}{1-z},\quad z(s)=\frac{s-5}{s+4},\quad
 \Re s(z)-\tfrac12=\frac{9(1-|z|^2)}{2|1-z|^2}.
\]
In particular $s(0)=5$ and the map is biholomorphic between the unit disk
and $\Re s>1/2$. Write $f(z)=H_F(s(z))=\sum_{m\ge0}v_m z^m$ locally
at zero. These $v_m$ are real. Introduce the ordinary Taylor coefficients
$u_k=\sum_{n\ge1} b_n(\log n)^k/(k!n^5)$ of $H_F(5-w)$. Near zero,
$5-s(z)=-9z/(1-z)$. Expanding $(1-z)^{-k}$ gives
\[
 v_0=u_0,\qquad
 v_m=\sum_{k=1}^m(-9)^k\binom{m-1}{k-1}u_k .
\]
For example, $v_1=-9u_1$, $v_2=-9u_1+81u_2$ and
$v_3=-9u_1+162u_2-729u_3$.

\paragraph{Certified finite approximation.}
The bound $|u_k-A_k/D_k|\le9(d+1)/N$ is uniform in $k$.
For completeness, comparing the right rectangle sum for $1/x$ on
$[1,n]$ with mesh $1/N$ to its integral, then allowing $1/N^2$ per
floor, gives $0\le\log n-J_n/N^2\le n/N$.
For $k\ge1$ the derivative of $t^k/k!$ there is bounded by $n$.
Also $|b_n|\le3(d+1)n$ and $(\log n)^k/k!\le n$. The first $N$
terms therefore cost at most $3(d+1)\sum n^{-2}/N\le6(d+1)/N$,
and the tail costs at most $3(d+1)\sum_{n>N}n^{-3}\le3(d+1)/N$.
For $k=0$ the prefix has zero error. The binomial theorem gives
$\sum_{k=1}^m9^k\binom{m-1}{k-1}=9\cdot10^{m-1}$. Clearing the
common denominator now gives
\[
              |v_m-T_{m,N}/D_m|\le E_m/N .
\]
At $N=3$ the RH packet gives approximations $-299/1944$, $383/1944$,
$-121/5832$ for the first three coefficients. These are small-cutoff
exact arithmetic examples; the certified errors here are large.

\paragraph{Forward implication.}
Under \eqref{eq:cayleysum}, partial summation gives
$H_F(s)=s\int_1^\infty S(x)x^{-s-1}\,dx$ for $\Re s>1/2$,
where $S(x)=\sum_{n\le x}b_n$. On $|z|=r<1$ we have
$|s|\le9/(1-r)$ and $\delta=\Re s-1/2\ge(9/4)(1-r)$.
Substituting $x=e^t$ gives the integral majorant
$100C|s|(\delta^{-1}+2\delta^{-2}+2\delta^{-3})$.
Thus $|f(z)|\le4500C(1-r)^{-4}$. At $r=m/(m+1)$, Cauchy's
estimate and $(1+1/m)^m<3$ give $|v_m|\le V_m$. Adding the finite
approximation error proves every assertion in \eqref{eq:cayleytest}.

\subsection*{Recovering the missing sign}
A strict one-sided estimate on the growth order is generally insufficient
for an arbitrary power series. Here the positive real interval is special:
$s(x)\ge5$ when $0\le x<1$. We prove the requisite positivity
argument directly and give finite translation bounds.

It is enough to assume the upper comparison at $N=E_m$. It gives
$v_m\le V_m+2$; the error $1$ occurs once in the test and once in passing
to the true coefficient. For the lower comparison replace $f$ by $-f$.
Let $K=13500C+2$, write $g=f$ in the upper case and $g=-f$ in the
lower case, and denote the coefficients of $g$ by $g_j$. Then
$c_j=K(j+1)^4-g_j\ge0$, including $j=0$ since $|u_0|\le6(d+1)$.
Put
\[
 G(z)=K\sum_{j\ge0}(j+1)^4z^j-g(z).
\]
The first term is rational on the disk, obtained by applying
$(1+z\,d/dz)^4$ to $(1-z)^{-1}$. The definition of $g$ by its
Dirichlet series supplies analytic neighborhoods along $[0,1)$.

We claim $c_m r^m\le G(r)$ for rational $0<r<1$. Subtract the finite
prefix $P(z)=\sum_{j=0}^m c_j z^j$. All Taylor coefficients of $G-P$
at $0$ are nonnegative. At a real center $x$, translating its Taylor
coefficients through a positive distance $h$ inside a disk of convergence
uses only nonnegative weights $\binom{j+k}{j}h^k$. Thus nonnegativity
passes along finitely many overlapping disks to $r$, proving
$G(r)-P(r)\ge0$. Since $P(r)\ge c_m r^m$, this is the claim.
This also explains the classical Pringsheim argument in this setting.

For an explicit mesh take $R=(1-r)/8$ and positive steps $h\le R/4$.
On $|z-x|\le R$, $0\le x\le r$, we have $\Re s(z)\ge4$. Indeed
the real part of $1/(1-z)$ is at least $1/(1-x+R)$ on this disk, and
$s=9/(1-z)-4$. Hence $|g(z)|\le3(d+1)\sum n^{-3}\le6(d+1)$
there. Since $(j+1)^4\le24\binom{j+4}{4}$,
\[
  G(r)\le 24K(1-r)^{-5}+6(d+1).
\]
At $r=m/(m+1)$ we obtain $c_m\le 3[24K(m+1)^5+6(d+1)]$ and hence
$|v_m|\le100K(m+1)^5$. The Taylor series now converges absolutely on
the whole disk. It agrees near zero with the meromorphic continuation
of $H_F(s(z))$ and by uniqueness cannot hide any pole there. The pole
comparison proves the converse for each family.

\paragraph{How many derivatives and bits does positivity require?}
This detail matters if induction is to be restricted to bounded,
decidable computations. All derivatives at a point need never be
propagated in a single induction invariant. Let $L(z)=G(z)-P(z)$.
A uniform integer bound $M\ge1$ on $|L|$ on the mesh disks is obtained
by rounding up
\[
 1+24K\bigl(1-r-R\bigr)^{-5}+6(d+1)
       +\sum_{j=0}^m \bigl(K(j+1)^4+54(d+1)10^j\bigr).
\]
Indeed $|u_k|\le6(d+1)$ follows by absolute summation, hence the
unconditional rough bound on $|v_j|$ used in this expression.
Let $w_j(x)=R^j L^{(j)}(x)/(j!M)$. Cauchy's estimate gives
$|w_j(x)|\le1$. Translation from $x$ to $x+h$ yields
\[
 w_j(x+h)=\sum_{k\ge0}\binom{j+k}{j}(h/R)^k w_{j+k}(x).
\]
The total weight is at most $2^{j+1}$ and the tail after $k=b$ is at most
$2^{j-b}$, using $\binom{j+k}{j}\le2^{j+k}$ and $h/R\le1/4$.
To prove $w_j(x+h)\ge-2^{-e}$ for every $0\le j\le J$, retain
$0\le k\le J+e+1$. It is enough to know
$w_i(x)\ge-2^{-(e+J+2)}$ for $i\le2J+e+1$. The weighted error
and the tail each cost at most $2^{-e-1}$. Start at the endpoint with
$(J,e)=(0,t)$ and work backwards by
\[
             (J,e)\longmapsto(2J+e+1,e+J+2).
\]
There are at most $\lceil32r/(1-r)\rceil$ steps. Every requested
entry at zero is nonnegative, so $L(r)\ge-M2^{-t}$ for arbitrary
$t$. Taking this as the definition of closed order on rational Cauchy
names gives $L(r)\ge0$. The mesh, order, precision, and sums are
primitive recursive in the rational data. This argument isolates the
finite propagation from the analytic certificates that justify Taylor
translation and its tails.

\paragraph{Formal scope.}
\texttt{CayleyPackets.lean} constructs the 18 closed object formulas.
\texttt{matrix\_reading} and \texttt{statement\_reading} identify their
integer meanings. \texttt{weightSum} proves the exact binomial amplification.
\texttt{translation\_step} proves the finite weighted error inequality over
rationals and \texttt{finite\_propagation} verifies the backwards budget
recursion. The object derivations \texttt{scheduled} and
\texttt{certified\_failure} identify, respectively, equivalent single-index
cube tests and least-radius positive counterexample certificates in iPRA.
The holomorphic identification, Cauchy estimates, and analytic re-expansion
on the mesh are human proofs. The Lean development supplies neither an
internal iPRA derivation from the prime-gap sentence to the one-sided tests
nor an internal derivation of the analytic limit identities.

\section{Multiplying polynomials on a Taylor disk}\label{sec:hardy-toeplitz}

A coefficient bound detects the radius of convergence. A different test asks what happens to the coefficients of every integer polynomial when it is multiplied by the Taylor series. Truncating the product gives a triangular Toeplitz matrix. The full test controls its largest singular value. A smaller test uses just its first column. Both admit explicit error allowances at every finite precision, including when the matrix is singular.

\paragraph{Analytic and arithmetic input.}
We use the prime-product construction of Section~\ref{sec:foundation} and the Taylor estimates in the companion report. For the monic polynomial $F$ of degree $d$ and height parameter $H$, let $B(n)$ be the integer part of the binary logarithm of its prime product, with $B(0)=0$. The three families use $F=X-1$, every $X^q-1$, or every monic integral $F$, respectively. Write $a_n=B(n)-B(n-1)$ and $b_n=a_n-2\mathbf1_{2\mid n}a_{n/2}$, where the last term is zero for odd $n$. Use the series
\[
 H_F(s)=\sum_{n\ge1} b_n n^{-s}\qquad(\Re s>1).
\]
Its continuation has a pole at every zero of the field zeta product with $1/2<\Re s<1$. Indeed the binary-log rounding contributes the Mellin transform of an error bounded by one, holomorphic for $\Re s>0$. The remaining logarithmic derivative has residues of one sign at zeros. The multiplier $1-2^{1-s}$ cancels the pole at $1$ and does not vanish in the interior of the strip. The bad-prime corrections to the field logarithmic derivatives are holomorphic for $\Re s>0$.

For $i\ge0$ put $c=4+2^i$, $Q=2c$, $P=(2c-1)c-1$ and $r=P/Q$. Write $C=1000(d+1)^4(H+2)^2$ and $U=1000cCQ^3$. Under the relevant zero hypothesis the prime-product estimate in the attached report implies $|\sum_{n\le x}b_n|\le 100C\sqrt{x}(1+\log x)^2$. Partial summation on $|s-c|\le r$ gives $|H_F(s)|\le U$. Here $|s|\le2c$, $\Re s\ge1/2+1/Q$ and
$\int_0^\infty(1+t)^2 e^{-t/Q}\,dt\le5Q^3$. This implication uses the attached explicit prime-product theorem as its analytic input.

Write $H_F(c-z)=\sum_{k\ge0}u_k z^k$ locally at zero. For $N\ge1$ use the integers
\[
 \begin{split}
 J(n,N)&=\sum_{j=1}^{N(n-1)}\left\lfloor\frac{N^2}{N+j}\right\rfloor,\qquad
 D_k=k!N^{2k}(N!)^c,\\
 A_k&=\sum_{n=1}^N b_n J(n,N)^k(N!/n)^c.
 \end{split}
\]
Then $|A_k/D_k-u_k|\le9(d+1)/N$ for every $k$. Indeed $u_k=\sum b_n(\log n)^k/(k!n^c)$, $|b_n|\le3(d+1)n$, and right Riemann sums followed by a floor error give $0\le\log n-J(n,N)/N^2\le n/N$. The derivative of $t^k/k!$ on $[0,\log n]$ is at most $n$ when $k\ge1$. Summing the resulting errors costs at most $3(d+1)\sum_{n\le N}n^{3-c}/N\le6(d+1)/N$. For the tail use $(\log n)^k/k!\le n$ and $c\ge5$. The $k=0$ case has zero rounding error.

\paragraph{Three finite tests.}
At length $m\ge1$ take a common denominator and an integer budget
\[
 L=Q^m m! N^{2m}(N!)^c,\qquad
 S=L\{NU+9m(d+1)(1+c)^m\}.
\]
The lower triangular Toeplitz matrix $M$ has subdiagonal entries
\[
 T_k=N P^k A_k Q^{m-k}\frac{m!}{k!}N^{2(m-k)}
 \qquad(0\le k<m).
\]
Thus $T_k/(NL)=r^k A_k/D_k$. Write $(Mx)_i=\sum_{j=0}^{i}T_{i-j}x_j$ and $e_0=(1,0,\ldots,0)$. Each of the following assertions quantifies over all allowed $F,i,m,N$. All vectors have integer coefficients and length $m$.
\begin{center}
\begin{tabular}{@{}ll@{}}
First column & $\sum_{k<m}T_k^2\le S^2$.\\[2pt]
Polynomial multiplication & $\|Mx\|_2^2\le S^2\|x\|_2^2$ for every $x$.\\[2pt]
Block quadratic form & $2y^{\mathsf T}Mx\le S^2\|x\|_2^2+\|y\|_2^2$ for every $x,y$.
\end{tabular}
\end{center}

\paragraph{Proof of the analytic equivalence.}
Assume the zero hypothesis. On the unit circle $f(z)=H_F(c-rz)$ has modulus at most $U$ and extends to a neighborhood of the closed disk. For $X(z)=\sum_{j<m}x_j z^j$, apply $\int_0^{2\pi}|fX|^2\,dt/(2\pi)\le U^2\int_0^{2\pi}|X|^2\,dt/(2\pi)$, evaluating at $z=e^{it}$. Orthogonality gives $\sum_{k<m}|[z^k](fX)|^2\le U^2\sum_{j<m}|x_j|^2$. This follows by subtracting the first $m$ Fourier coefficients and integrating the resulting nonnegative square. In these first $m$ coefficients replacing $r^ku_k$ by $T_k/(NL)$ introduces a Toeplitz matrix whose entries have absolute value at most $\delta=9(d+1)(1+c)^m/N$. Its operator norm is at most $m\delta$. Indeed it is a sum of $m$ scalar multiples of truncated shifts, each summand of norm at most $\delta$. Thus $\|Mx\|_2\le NL(U+m\delta)\|x\|_2=S\|x\|_2$.

Testing $x=e_0$ gives the first-column condition. Conversely, assume only that condition. For each fixed $i,k$ take $m=k+1$ and let $N$ tend to infinity after dividing by $(NL)^2$. This gives $|r^k u_k|\le U$. The geometric majorant gives convergence throughout $|s-c|<r$. Uniqueness of meromorphic continuation excludes poles there. For a fixed $x+iy$ with $x>1/2$, put $a_c=1/2+1/(2c)$. Membership in the open disk is equivalent to $2c(x-a_c)>x^2+y^2-a_c^2$, eventually true along $c=4+2^i$. Every off-line zero in the right half of the strip would therefore produce a pole in one of these disks. Reflection and the identification of the three polynomial families with their zeta or character products finish the converse. The passage through analytic continuation is a standard-model argument.

\paragraph{A direct iPRA equivalence.}
The second and third tests are interderivable by finite arithmetic without the analytic bridge. For any integer $x,y$, the right side of the third inequality minus its left side is
\[
   S^2\|x\|_2^2-\|Mx\|_2^2+\sum_{i<m}(y_i-(Mx)_i)^2.
\]
The final sum is nonnegative. Conversely choose $y=Mx$. Integer entries are differences of naturals. For a concrete radix representation, store positive and negative parts of $x_i,y_i$ in four consecutive digits of a word, in radix $B\ge2$. To synthesize the witness, choose $B=2+\sum_{i<m}(|x_i|+|(Mx)_i|)$, form the digit word, and add $B^{4m}$. This high unused digit ensures the word bounds the radix offset. Digit extraction follows by division with remainder. Its induction invariant states the corresponding prefix remainder and quotient identities, and the square identity follows by summing the polynomial identity one term at a time. These are quantifier-free inductions on bounded accumulators. Substitution of the resulting primitive-recursive word witnesses the reverse implication. At fixed $m$ the block matrix is $\left(\begin{smallmatrix}S^2 I&-M^{\mathsf T}\\-M&I\end{smallmatrix}\right)$.

\paragraph{Recovering the archive's coefficient test inside arithmetic.}
There is also a direct finitary implication to the prior lacunary Taylor test. It asks for
\[
  r^k |A_{k,N}/D_{k,N}|\le U+Er^k/N,\qquad E=9(d+1).
\]
The numerator and denominator subscripts now indicate the cutoff explicitly. At first sight the first-column bound only gives a larger slack of order $m(1+c)^m/N$. The following finite argument recovers the stated allowance exactly.

Set $L_N(n)=J(n,N)/N^2$ and $R_N(n)=\sum_{j=1}^{N(n-1)}1/(N+j)$. If $V$ is a multiple of $N$, monotonicity of the right and left rectangle sums gives
\[
  0\le R_V(n)-R_N(n)\le (1-1/n)/N,\qquad
  0\le R_V(n)-L_N(n)\le n/N.
\]
These are finite comparisons: group the finer rectangles belonging to each coarse one, telescope the left-minus-right sum, and allow at most $(n-1)/N$ for the floors. A second finite bound replaces the exponential function. Write $t_j=1/(V+j)$ for $1\le j\le V(n-1)$. The multinomial expansion, followed by the finite geometric sum, gives for all $k\ge0$
\[
 \frac{R_V(n)^k}{k!}\le \prod_j\sum_{e=0}^k t_j^e
 \le \prod_j(1-t_j)^{-1}=n.
\]
Indeed the terms with total exponent $k$ acquire factors $1/e!$ in the multinomial expansion. The last product telescopes, also for $n=1$ with an empty product. For $0\le u\le t\le R_V(n)$ and $k\ge1$, factor $t^k-u^k$ as $(t-u)\sum_{j<k}t^{k-1-j}u^j$. This bounds its quotient by $k!$ by $n(t-u)$.

Put $v_N=A_{k,N}/D_{k,N}$, keeping $c,k$ fixed. To compare $v_N$ with $v_K$ take $V=NK$ and apply the preceding inequalities for each $n\le\min(N,K)$. Replacing $L_N,L_K$ by $R_V$ in the common terms costs at most
\[
  3(d+1)(N^{-1}+K^{-1})\sum_{n\le\min(N,K)}n^{3-c}.
\]
This uses the elementary prime-product jump bound $|b_n|\le3(d+1)n$. For $k=0$ the common terms agree exactly. The unmatched terms cost at most $3(d+1)/\min(N,K)$, by $L_N(n)^k/k!\le n$ and its counterpart at $K$. Both estimates follow from $c\ge5$ and $1/n^2\le1/(n-1)-1/n$ for $n\ge2$. Therefore
\[
  |v_N-v_K|\le E(1/N+1/K).
\]
All the grids, multi-indices, and products just used have explicit primitive-recursive bounds. The finite sum comparisons, telescoping, multinomial identity, and digit and degree bounds admit quantifier-free induction on their finite accumulators after clearing denominators.

Take $m=k+1$ in the first-column test at precision $K$. One term of its nonnegative sum of squares gives $r^k|v_K|\le U+Em(1+c)^m/K$. Combining with the finite comparison yields
\[
 r^k|v_N|\le U+Er^k/N+\frac{E\{r^k+m(1+c)^m\}}{K}
 \qquad(K\ge1).
\]
Multiply by positive common denominators fixed by $i,k,N$ and then by $K$. This has the form $KA\le KB+J$ for every positive integer $K$. If $A>B$, the single substitution $K=J+1$ contradicts it. This proves the prior coefficient test with precisely its old allowance. In particular the first-column implication here uses neither complex analysis nor an infinitary limit rule. Lean checks the finite power increment and slack removal in \texttt{monomial\_increment} and \texttt{remove\_allowance}; the concrete rational grid and prime-packet transfer are proved on paper.

\paragraph{A determinant criterion with strict comparisons.}
There is a fourth test with no vector inputs. Set $G=S^2 I-M^{\mathsf T}M$. For every $k\ge0$ require every \emph{leading} principal determinant of $I+kG$ to be strictly positive. The shift permits a singular $G$. For example, $M=5I$ and $S=5$ gives $G=0$, with every shifted determinant equal to one. The first column and full contraction can differ for a single packet. With $m=2$, $T_0=3$, $T_1=4$ and $S=5$ the first-column inequality is an equality, but $x=(1,1)$ gives $\|Mx\|_2^2=58>50=S^2\|x\|_2^2$. The first-column converse above uses all precisions and lengths.

To prove the determinant equivalence in finite arithmetic, first assume $x^{\mathsf T}Gx\ge0$ for integral $x$. Clearing denominators gives the same result for rational $x$. Hence $I+kG$ is strictly positive on each nonzero rational vector. For a symmetric rational matrix strict positivity on those vectors is equivalent to positivity of its leading determinants. Indeed the first pivot $a$ is positive and completing the square gives a first term $a(t+a^{-1}b^{\mathsf T}y)^2$ and residual matrix $D-a^{-1}bb^{\mathsf T}$. Strict positivity passes to the residual by taking $t=-a^{-1}b^{\mathsf T}y$. Every leading determinant that includes the first row is $a$ times the corresponding residual determinant, by row elimination. This also proves the converse inductively when all the leading determinants are positive. For the remaining implication fix integral $x$. If $x^{\mathsf T}Gx<0$, the integer is at most $-1$ and $x\ne0$. With $k=1+\|x\|_2^2$ one gets $x^{\mathsf T}(I+kG)x<0$, contradicting positivity of the leading determinants. The test checks all shifts without a bound on the size of the possible vector.

Here the finite elimination proof can be used with quantifier-free induction. Eliminate along successive positive pivots, keeping the rational matrices and unit triangular change of variables as coded arrays. The stage invariant consists of the factorization identity and the leading-determinant product identities, both bounded entrywise tests. The determinant identity for a row addition follows by pairing the equal-row permutation terms of the Leibniz expansion. A failed pivot, after the previous positive pivots, supplies a rational test vector by back substitution, and the product of its denominators clears it to an integral vector. Thus an induction on a formula quantifying over arbitrary vectors at every recursion step is unnecessary. All rational operations can be performed by signed numerators and positive denominators. The compiled test evaluates determinants by their finite signed permutation expansion.

\paragraph{Formal status.}
\texttt{HardyToeplitz.lean} builds the twelve closed sentences and proves object-language equivalence of each to its scheduled finite-cube versions. The theorem \texttt{statement\_check} identifies the first nine sentences with explicit inequalities in sums of integer digits and signed coefficient components. The integer lemmas \texttt{polarization\_iff}, \texttt{contraction\_iff\_block}, and \texttt{nonneg\_iff\_all\_shifts} check completion of squares and the integer shift argument. The paper supplies the internal arithmetic proof for the concrete probes and elimination, including the coding of witnesses. The Lean file constructs the determinant expansion directly as a primitive-recursive program. Its interpretation as a determinant and the Sylvester argument are proved here on paper. The analytic bounds, Fourier argument, and converse through pole detection are standard-model arguments.

\section{Robin and Nicolas with exact rational witnesses}\label{sec:named}
Let $T(n)=e^\gamma\log\log n$, where $\gamma$ is Euler's constant. Two classical criteria for ordinary RH compare this number with rational numbers in opposite directions. With $P_y=\prod_{p\le y}p$ they are
\[
 \frac{\sigma(n)}n<T(n)\quad(n\ge5041),\qquad
 T(P_y)<\frac{P_y}{\varphi(P_y)}\quad(y\ge3).
\]
Here $\sigma(n)=\sum_{a\mid n}a$ and $\varphi(P_y)=\prod_{p\le y}(p-1)$. Starting at $y=3$ means that the first primorial considered is six. A composite cutoff repeats a primorial and cannot change the assertion. Robin's lower cutoff is on $n$, including $5041$.

We construct rational brackets $T_m^-(n)\le T(n)\le T_m^+(n)$ with both ends tending to $T(n)$ for every allowed input. This supports two different assertions for each comparison. A strict comparison gives a finite successful precision as evidence. A weak comparison checks every rational enclosure for a violation. The quantifier order is part of the formulation.

\paragraph{Logarithms from rational rectangles.}
For rational $x\ge1$ and $m\ge1$ use the equal-width right and left sums for $1/t$:
\[
 \ell^-_m(x)=\sum_{j=1}^{m}\frac{x-1}{m+j(x-1)},\qquad
 \ell^+_m(x)=\sum_{j=0}^{m-1}\frac{x-1}{m+j(x-1)}.
\]
They enclose $\log x$ since the integrand is decreasing. Their difference is $(x-1)(1-1/x)/m$. At $x=1$ both vanish. Evaluate the first logarithm of an integer through binary range reduction. For $h=\lfloor\log_2 n\rfloor$ set
\[
 b_m^{\pm}(n)=h\ell_m^{\pm}(2)+\ell_m^{\pm}(n/2^h),\quad
 c_m^{\pm}(n)=\ell_m^{\pm}(\max(1,b_m^{\pm}(n))).
\]
All inputs used here have $\log n>1$. Monotonicity of the actual logarithm, followed by the rectangle inequality, therefore gives
$c_m^-(n)\le\log\log n\le c_m^+(n)$. For fixed $n$, the $b_m^{\pm}$ converge to $\log n$ and stay in a bounded interval. The rectangle error on that interval is $O_n(1/m)$, so both $c_m^{\pm}$ converge to $\log\log n$. At the code level, for $x=a/b$, set $e=(a-b)_+$, $w=mb$. The common denominator $\prod_{j=0}^{m}(w+je)$ is positive and each selected numerator contribution is $e$ times an exact quotient. Thus every division in the evaluation has a specified nonzero denominator.

\paragraph{Bounding Euler's constant.}
Let $H_m=\sum_{i=1}^m 1/i$ and $r=2m$. Use
\[
 g^-_m=\frac{(1+H_m/r)^r}{m+1},\qquad
 g^+_m=\frac{1}{m(1-H_m/r)^r},\qquad
 T_m^\pm(n)=g_m^\pm c_m^\pm(n).
\]
The increasing sequence $H_m-\log(m+1)$ and decreasing sequence $H_m-\log m$ have common limit $\gamma$. This follows by comparison with the integral of $1/t$ on each unit interval. The binomial factors enclose $e^{H_m}$ by $\log(1+x)\le x\le-\log(1-x)$ for $0\le x<1$. Here $H_m/r\le1/2$ and $H_m\le1+\log m$. The logarithmic errors incurred in either binomial factor are at most $r(H_m/r)^2$ by integrating the errors of the two logarithmic derivatives on $[0,H_m/r]$. This tends to zero with $m$. Hence $g^-_m\le e^\gamma\le g^+_m$, with convergence at both ends. The products defining $T_m^\pm$ preserve the inequalities since all factors are nonnegative.

\begin{center}\small
\begin{tabular}{@{}lll@{}}\toprule
Criterion & Strict certificate for each input & Weak order at every input and precision\\\midrule
Robin & $\exists m\ge1:\ \sigma(n)/n<T^-_m(n)$ &
 $\sigma(n)/n\le T^+_m(n)$\\
Nicolas, $P=P_y$ & $\exists m\ge1:\ T^+_m(P)<P/\varphi(P)$ &
 $T^-_m(P)\le P/\varphi(P)$\\\bottomrule
\end{tabular}
\end{center}
\begin{theorem}
Each of the four sentences in the table characterizes ordinary RH in the classical standard model. The strict rows are $\Pi^0_2$ and the weak rows are $\Pi^0_1$ in the PR language.
\end{theorem}
\begin{proof}
Rational comparison is exact after multiplying by positive denominators. A strictly positive real gap is eventually certified by the indicated endpoint, since both endpoints converge. A successful strict finite comparison implies the original strict inequality. A weak test passes at all precisions exactly when the limiting weak comparison holds. For Robin's theorem the failure of RH produces infinitely many strict reverse inequalities. For Nicolas's theorem, off RH the signed normalized deviation below is unbounded in both directions on primorials, so there are strictly negative deviations. Thus allowing equality in the weak versions preserves the classical RH characterization. This step uses the oscillation results: enclosure convergence alone would not turn a weak inequality into a strict one. See Robin (1984), as stated with its strengthened converse in \href{https://arxiv.org/pdf/math/0008177}{Lagarias, Theorems 3.1--3.2}, and \href{https://arxiv.org/pdf/1202.0729}{Nicolas, Theorem 1.1 and its converse discussion}.
\end{proof}

\paragraph{A bounded-error criterion for primorials.}
Write $R=P/\varphi(P)$ and $c(P)=(R-T(P))\sqrt{\log P}$. Under RH, Nicolas proves $0<c(P)<5$ for the primorials considered here (in fact the global maximum in his bound is about $4.063$). If RH fails, the primorial values are unbounded above and below. Thus each of $c(P)\le5$, $c(P)\ge-5$, and $|c(P)|\le5$ for all these primorials characterizes RH. The three corresponding universal rational tests use
\[
  b^-_m(P)\,(R-T^+_m(P))_+^2\le25,\qquad
  b^-_m(P)\,(T^-_m(P)-R)_+^2\le25.
\]
Use respectively the first, the second, or their conjunction, with $m\ge1$. Each left side approaches its exact signed-part square from below, by the enclosures above and continuity of positive part. Therefore the universal test admits equality and introduces no square root or real constant. This use of a fixed $5$ is justified by Nicolas's quantitative theorem; an arbitrary positive number would change the claim.

\paragraph{Lean and examples.}
\code{RobinNicolas.lean} constructs all seven sentences using rational numerators and denominators, the primitive recursive divisor sum, a bounded prime product, and the factorial numerator for $H_m$. The theorems \code{strict_rational_reading}, \code{weak_rational_reading} and \code{scaled_rational_reading} interpret those formulas as inequalities of exact rationals, with positivity of the denominators checked. For each strict comparison it also proves \code{strict_iff_first} in object iPRA. Given a successful precision, bounded search through that precision produces its least successful one; the verified invariant keeps the first positive margin. This argument preserves the order $\forall n\exists m$, without choosing a function of $n$. With the actual enclosures above the least successful $m$ at $n=5041$ in the Robin test is 2. For $P=6,30,210,2310$ in the Nicolas test the values are $3,7,12,21$. They are exact rational comparisons. These finite cases illustrate the meaning of a witness.

\section{Ordinary RH as a finite least-squares certificate}\label{sec:residue-space}

Let $R_k(n)=n\bmod k$ for $k\ge2$, $n\ge1$. Our test asks how well a
finite rational linear combination of these periodic sequences approximates
the constant sequence $1$. For example, $R_2$ agrees with $1$ on odd
integers and vanishes on even integers. The weight of an error at $n$ is
$w_n=1/[n(n+1)]$. Its exact tail $\sum_{n>J}w_n=1/(J+1)$ makes
a finite certificate possible.

\paragraph{The analytic input.}
Bagchi's Theorem~1 states that ordinary RH is equivalent to $1$ belonging
to the closed span of $\{n/k\}$ in the complex sequence space
$\ell^2(w)$, where the braces mean fractional part. Multiplying each
generator by its index gives $R_k$ and leaves its span unchanged. Real
parts of coefficients suffice since the target is real and this operation
reduces every residual magnitude. Rational coefficients suffice because changing a finite list of coefficients by $\eta_k$ changes the norm by at
most $\sum_k(k-1)|\eta_k|$, since the weights sum to one. The source is
B.~Bagchi, \emph{On Nyman, Beurling and B\'aez-Duarte's Hilbert space
reformulation of the Riemann hypothesis}, Proc. Indian Acad. Sci.
(Math. Sci.) 116 (2006), 137--146,
\href{https://arxiv.org/abs/math/0607733}{arXiv:math/0607733}, Theorem~1.
This use of Bagchi is a classical standard-model identification.

\paragraph{Three finite-arithmetic readings of approximation.}
Choose a denominator $D\ge1$ and nonnegative integers $u_k,v_k$ on
$2\le k\le K+1$, where $K$ may be zero. The approximant has coefficients
$(u_k-v_k)/D$. For any of the zero-one targets $f$ below, put
\[
 g(n)=\left|Df(n)-\sum_{k=2}^{K+1}(u_k-v_k)R_k(n)\right|,\qquad
 B=D+\sum_{k=2}^{K+1}(u_k+v_k)(k-1),\qquad
 E_N=\sum_{n=1}^N\frac{g(n)^2}{n(n+1)}.
\]
Every residual is bounded by $B$, irrespective of its sign.
For each integer precision $t\ge1$, ask for one approximant satisfying
one of the following tests.
\begin{enumerate}
\item Every prefix satisfies $tE_N\le D^2$, including $N=0$.
\item Some $J\ge0$ satisfies $t(E_J+B^2/(J+1))\le D^2$.
\item The same finite inequality holds at $J=2tB^2$.
\end{enumerate}
Tests 2 and 3 are $\Pi^0_2$ presentations with a primitive-recursive
matrix. Test 1 has an additional universal block. Quantifying over
arbitrary precisions allows the approximant to improve with $t$. A
computable selector for those approximants is a separate assertion.

The finite inequality in tests 2 and 3 controls every prefix. Indeed,
$E_N$ increases, whereas $E_N+B^2/(N+1)$ decreases, because the change in
the latter after adding index $N+1$ is
\[
       \frac{g(N+1)^2-B^2}{(N+1)(N+2)}\le0.
\]
For $N\le J$ use the first monotonicity, and for $N\ge J$ use the
second. Conversely, use the witness for test 1 at precision $2t$.
At $J=2tB^2$ the prefix contributes at most $D^2/(2t)$ and the
allowance contributes at most $1/(2t)\le D^2/(2t)$. This gives
test 3; supplying this $J$ also gives test 2. Classically, bounding all
prefixes is exactly bounding the squared norm of the residual. Approximants
of arbitrarily small norm and witnesses for every integer $t$ therefore
express the same closure assertion for $f=1$.

\paragraph{Equivalence of the constant, step and coordinate targets.}
The same three tests may instead require approximation to every
$h_m(n)=[n\ge m]$, or to every $e_m(n)=[n=m]$, with $m\ge1$ quantified
before the approximant. All nine resulting sentences are constructively
equivalent by finite witness transformations. Here are details that avoid
selecting a growing list of existential witnesses.

Extend any sequence by zero at index zero, and put
$(T_m a)(n)=a(\lfloor n/m\rfloor)$. For $j\ge1$ the weight of the complete
block $mj\le n<m(j+1)$ is $w_j/m$, by telescoping. A partial last block
has at most that weight. Thus the squared error of a dilation through any
cutoff $N$ is at most $1/m$ times the original squared error through
$\lfloor N/m\rfloor$. Dilation sends a residue vector to a rational linear
combination of residue vectors by the identity
\[
   T_mR_k=(R_{mk}-R_m)/m,\qquad R_1=0.
\]
This follows by writing $n=mq+b$, $0\le b<m$, and then
$q=ka+c$, $0\le c<k$. The remainder of $n$ modulo $mk$ is $mc+b$.
A denominator for the new vector is $mD$.

Since $T_m1=h_m$, one witness for $1$ gives one for any $h_m$ at the
same precision. Conversely $h_1=1$. For $e_m=h_m-h_{m+1}$, take the
difference of two dilations of the \emph{same} approximation to $1$ at
precision $4t$. Using $(x-y)^2\le2x^2+2y^2$ gives squared error at most
$1/t$ on each prefix. Conversely, take a witness approximating $e_1$.
For $L=4t$,
\[
   \sum_{i=0}^{L-1} T_{2^i}e_1=[1\le n<2^L].
\]
This union comprises successive disjoint intervals; its squared
distance from $1$ on every prefix is at most $2^{-L}$. Use one witness
for $e_1$ at precision $4tL^2$ and sum those $L$ dilations. By the finite
Cauchy inequality $(\sum_i x_i)^2\le L\sum_i x_i^2$, the resulting
squared error from the displayed indicator is at most $1/(4t)$.
Since $2^L\ge4t$, the squared error from $1$ is at most
$2/(4t)+2/2^L\le1/t$. This also proves that the single target $e_1$
suffices. Multiplying denominators clears every coefficient in these
constructions. Every loop length, new modulus, denominator and coefficient
word is a primitive-recursive function of the given witnesses and
precision. The finite Cauchy inequality follows by summing
$(x_i-x_j)^2\ge0$ over $i,j$. Thus the needed inductions are inductions
on finite sums after denominator clearing, with quantifier-free
invariants. The argument supplies an iPRA proof on paper of the links
among the residue tests; the link to the prime-product sentences through
ordinary RH uses the external analytic theorem.

\paragraph{Encoding and verification.}
In \texttt{ResidueSpace.lean}, the two words encode the $K$ nonnegative
digits in a common base $H+2$. Every finite signed list admits this coding
by taking a base larger than its greatest absolute numerator. The file
constructs nine closed \texttt{Fm PTm 0} sentences. Let
$P_N=\prod_{j=1}^N j(j+1)$. The encoded numerator $S_N$ is a sum over
$n\le N$ of $g(n)^2$ times the product of all factors in $P_N$ except
the $n$th. This includes $P_0=1$, $S_0=0$.
The certificate is exactly
\[
 t\big((J+1)S_J+B^2P_J\big)\le D^2(J+1)P_J.
\]
The Lean arithmetic lemmas formalize the residual bound, denominator
clearing, monotone prefix and tail estimates, the two precision
conversions and the mixed-radix identity for dilation. The theorems
\texttt{closed\_fixed\_iff\_prefix} and
\texttt{closed\_cert\_iff\_prefix} apply directly to the standard
interpretations of the compiled sentences for each of the three targets. The preceding
finite block and target-conversion proofs supply the rest of the argument.
For a small example, $D=1$, $u_2=1$, $v_2=0$ and $f=1$ give $B=2$.
At $J=24$ the upper estimate equals $11973461273/26771144400<1/2$, so this
single row certifies precision $t=2$. The script
\texttt{check\_residue\_space.py} independently checks this example, the
denominator clearing, the one-step inequalities and the dilation identities
using exact rational arithmetic.

\subsection{Eliminating the approximating vector by two integer determinants}
The preceding certificates contain an arbitrary list of rational coefficients.
Two natural numbers, a matrix dimension $K$ and an observation cutoff $J$, suffice as witnesses.

For each of the same targets $f$ let the columns be
$h_i=R_{i+2}$ for $0\le i<K$, and $h_K=f$. Put
$H_i=i+1$ for $i<K$ and $H_K=1$, so $|h_i(n)|\le H_i$. Form the
$(K+1)\times(K+1)$ integer matrix
\[
 G_{ij}=(J+1)\sum_{n=1}^J h_i(n)h_j(n)
              \prod_{\substack{1\le a\le J\\a\ne n}}a(a+1)
       +[i=j]P_J(K+1)H_i^2.
\]
Let $A$ be its leading $K\times K$ submatrix and
$S=(J+1)P_J$. At $K=0$, $\det A=1$. The new test is
\[
 \boxed{\qquad \forall t\ge1\ \exists K,J\ge0\qquad
                   t\det G\le S\det A.\qquad}
\]
For a variable target index $m$, quantify $m$ before $t$. Every entry
of both matrices is a nonnegative integer. Their sizes, their entries
and both signed Leibniz sums are primitive recursive in $m,t,K,J$.

To prove the test, take a candidate coefficient vector $x\in\mathbb Q^K$,
write $e_x=f-\sum_i x_ih_i$, and put
\[
 Q(x)=\sum_{n=1}^J w_ne_x(n)^2+
            \frac{K+1}{J+1}\Big(1+\sum_{i<K}H_i^2x_i^2\Big).
\]
The second term bounds the entire squared error past $J$ because
$e_x(n)^2\le(K+1)(1+\sum H_i^2 x_i^2)$ by finite Cauchy, and the
total remaining weight is $1/(J+1)$. Thus $t Q(x)\le1$ gives
every-prefix approximation at precision $t$.
Conversely obtain one rational approximant $x$ with the every-prefix
bound at $2t$ and write $x_i=a_i/D$. Choose
\[
   J=2t(K+1)\Big(1+\sum H_i^2 a_i^2\Big).
\]
Since $D\ge1$, the tail allowance in $Q$ is at most $1/(2t)$.
It follows that $tQ(x)\le1$. This proves equivalence before eliminating
the coefficient vector.

Write $G=\left(\begin{smallmatrix}A&b\\ b^T&d\end{smallmatrix}\right)$.
Expansion of the integer matrix gives
$S Q(x)=x^TAx-2b^Tx+d$.
Both $G$ and $A$ are positive definite because each is a sum of nonnegative
outer products and a positive diagonal matrix. Their determinants are
positive over the rationals. To see this, induct on dimension. The first diagonal entry is a positive pivot. Subtracting
its multiples from lower rows yields its Schur complement. A nonzero
vector $y$ for that complement extends to
$(-a_{11}^{-1}\sum_{i>1}a_{1i}y_i,y)$, so its quadratic value is
positive. The determinant equals the pivot times the determinant of
the smaller positive definite matrix. Include the empty determinant
as the induction base. Hence $z=A^{-1}b$ exists over $\mathbb Q$, and
\[
 S Q(x)=(x-z)^TA(x-z)+d-b^Tz,\qquad
 \det G=(\det A)(d-b^T z).
\]
The second identity follows by subtracting $z^T$ times the first
$K$ rows from the last row. The minimum is attained at $x=z$, proving
the displayed determinant test in both directions. If the test passes,
Cramer's rule recovers integer numerators and common denominator
$\det A$ for the witness. This construction uses no minimization over
reals. In iPRA, common-denominator integer arithmetic suffices for the
elimination. Each stage of rational row reduction is a primitive
recursion; verifying the finite products, row identities and positivity
uses quantifier-free induction. For the positivity argument one may
carry the original diagonal lower bound through elimination because substitution of the explicit extension of $y$ retains at least the
sum of the original lower-bound terms on the surviving coordinates.
After clearing all positive denominators this is a coefficient identity
and a sum of squares. Thus it requires no induction on a predicate
quantifying over all rational vectors.

\texttt{ResidueGram.sentence} builds the three closed formulas.
\texttt{ResidueGram.eliminate} proves the finite rational equivalence
between the coefficient search and the determinant comparison for a
symmetric invertible core with nonnegative quadratic form and positive
determinant. The lemmas \texttt{ridge\_square} and \texttt{ridge\_pos}
also verify the sum-of-squares expansion and strict positivity of any
weighted design matrix with a positive diagonal penalty. Identifying the
compiled integer entries with this matrix and proving its determinant
positive are carried out above on paper. In the executable checker, the empty matrix and
cutoff are included explicitly.

\section{Counting down a failed packet inside the proof calculus}
The prime-product sentences test whether a nonnegative integer excess
$z$ vanishes. One may instead test for the last positive value of a
countdown. Define $p(0)=0$, $p(s+1)=s$ and
\[
 D(0,z)=z,\qquad D(k+1,z)=p(D(k,z)).
\]
The three following conditions are equivalent by derivations in
first-order iPRA, uniformly in the term $z$:
\[
 z=0,\qquad D(p(z),z)=0,\qquad
             \forall k\;\neg(D(k,z)=1).
\]
In the third condition, a failing packet has an explicit witnessing time.
For example, at excess $4$ the observations are $4,3,2,1,0,0,\ldots$,
and $p(4)=3$ selects the unique occurrence of $1$. A zero excess gives
the constant zero sequence.

Here is a proof using only the permitted induction scheme. Quantifier-free
induction on $k$ first gives $D(k,0)=0$. With a free parameter $z$, another
such induction gives
$D(k+1,z+1)=D(k,z)$. At $k=0$ the identity says $p(z+1)=z$.
For the successor step apply $p$ to both sides of the induction
hypothesis. Consequently $D(k,k+1)=1$ by induction on $k$, since its
initial value is $D(0,1)=1$ and the preceding cancellation identity
reduces each successor to its predecessor. Finally iPRA proves the
quantifier-free disjunction
\[
          z=0\quad\lor\quad D(p(z),z)=1 .
\]
At zero take the left branch. At a successor $z=s+1$, rewrite
$p(s+1)=s$ and take the right branch using $D(s,s+1)=1$.
From this disjunction either alternative test implies $z=0$: the other
branch would give $1=0$ for the computed test, or produce the prohibited
event at $k=p(z)$ for the universal test. In the opposite direction
substitute $z=0$ and use $D(k,0)=0$.

The file \texttt{CountdownTests.lean} uses the repository's
\texttt{predR} and \texttt{reversedSub} for these two primitive-recursive
operations. The theorem \texttt{equivalence} in that module lifts the pointwise
biconditional through any initial universal prefix ending in an equality
to zero. At each binder, specialize the assumed universal formula at a
fresh variable, apply the pointwise implication, then universally
introduce the variable again. Its specializations
\texttt{CountdownTests.rh}, \texttt{.grh} and \texttt{.erh} therefore
give six closed derivations of type \texttt{Prf PRAAx []}. The same
argument puts the computed test at the original quantifier complexity.
The universal test adds one universal variable and retains a
quantifier-free matrix. Each statement concerns the arithmetic packet
itself and uses no analytic bridge.

\subsection{Removing a fixed additive error by one finite multiplier}

A frequent final step in a rational estimate is to remove an additive
rounding allowance $J$ that is independent of precision. Suppose the
discrepancy has already been cleared of denominators and its positive part
is $z\in\mathbb N$. The following are interderivable inside first-order
iPRA, for each $J$:
\[
 z=0,\qquad (J+1)z\le J,\qquad
                 \forall k\ k z\le J .
\]
One can even allow $J$ to be an existential witness. Thus
$z=0\leftrightarrow\exists J\,\forall k\,kz\le J$. The order of the
quantifiers matters. For $z=1$ the variant
$\forall k\,\exists J\,kz\le J$ is satisfied by $J=k$.
In the motivating comparison $ka\le kb+J$ for naturals $a,b$,
ordinary finite arithmetic rewrites the discrepancy as
$z=a\mathbin{\dot-} b$. Its positive part obeys $kz\le J$, hence $a\le b$.
The Lean derivation below treats the excess $z$ as its input.

\paragraph{An object proof of the separating multiplier.}
Use the countdown $D(j,x)$ from the preceding argument. Since
$D(j,x)=x\mathbin{\dot-}j$, the order test $x\le j$ is the quantifier-free
equation $D(j,x)=0$. Define $M(0,k)=0$ and
$M(z+1,k)=k+M(z,k)$. This computes $kz$ with suitably ordered summands.
First prove
\[
              D(j,(j+1)+y)=y+1
\]
by quantifier-free induction on $j$. At $j=0$ it is the defining
zero-step equation. For the successor, rewrite the input as
$1+((j+1)+y)$ and apply $D(j+1,x+1)=D(j,x)$.
It follows that at a positive excess,
\[
           D(J,M(z+1,J+1))=M(z,J+1)+1.
\]
The right side is a successor, so it cannot equal zero. For $z=0$ all
the multiples are zero. This proves the finite test at $J+1$ and, by
specializing a universal hypothesis, the universal test. An existential
budget is eliminated by this same derivation with the supplied $J$.
In the forward direction zero itself can serve as a budget.
Every induction in this proof has a quantifier-free invariant.

\texttt{PrecisionBounds.finite\_recovery},
\texttt{fixed\_budget} and \texttt{bounded\_iff} are object derivations.
Applying the last theorem to the excess in each of the three primary
sentences allows a separate budget for each packet and gives three more
closed biconditionals of type \texttt{Prf PRAAx []}. In particular,
no uniform budget across different number fields, conductors or cutoffs
is assumed. The multiplier $J+1$ eliminates the budget for its own packet.

\section{Verification and the earlier criteria}\label{sec:status}
\paragraph{Using the archive.}
The Lean project builds against Lean 4.32.0 with its Mathlib revision pinned by the manifest. The entry point is \code{RiemannIPRA.lean}. The shell commands in the README rebuild the sources and re-run the arithmetic examples. The files under \code{verification/} include the build result, theorem-type and axiom audit, and exact example outputs. Kernel axiom reports can list \code{Classical.choice} for Lean's ambient matrix or finite-sum theorems. An object derivation remains a term built from ProveIt's scoped intuitionistic deduction rules. In particular, kernel axiom output alone does not distinguish the two forms of proof.

\paragraph{Locating the proofs.}
The notation in the last column specifies how far each correspondence has been checked. A standard-model reading interprets the actual compiled closed formula in the natural numbers.
\begin{center}\small
\begin{tabular}{@{}>{\raggedright\arraybackslash}p{.23\linewidth}>{\raggedright\arraybackslash}p{.35\linewidth}>{\raggedright\arraybackslash}p{.35\linewidth}@{}}\toprule
Criteria & Kernel-checked part & Analytic or finite argument in the paper\\\midrule
One-sided sampling & Object derivations for splitting signs and selecting samples. Standard-model reading and finite interpolation. & Landau's argument from the sparse upper bound back to zero-freeness.\\[2pt]
Cayley coefficients & Compiled integer readings, finite propagation, and object transformations of universal tests. & Approximation error, coefficient growth and analytic continuation.\\[2pt]
Taylor convolution & Integer readings, finite quadratic and determinant identities, and object changes of schedule. & Uniform analytic estimates, transfer between precisions, and the converse through the first column.\\[2pt]
Robin and Nicolas & Rational comparisons in the standard model. Object derivations for the first successful precision. & Convergence of the enclosures and the classical oscillation theorems.\\[2pt]
Residues and Gram matrices & Closed readings and conversions between cutoff tests. Rational elimination and positivity lemmas. & Changes of target, identification of the integer determinants, and the link to Hilbert-space density.\\[2pt]
Countdowns and budgets & All displayed pointwise equivalences and their specializations to the primary hypotheses are object derivations. & Only the interpretation of the primary hypotheses invokes analysis.\\\bottomrule
\end{tabular}
\end{center}

\paragraph{What the attachments contribute.}
The complete cumulative mathematical report from the third attachment is included as \code{background/report.pdf}, with its TeX sources beside it. The new development reuses its actual Lean modules. The first attachment introduced the polynomial prime-power tests, the second added the arithmetic equivalences, and the third added the splitting-prime, Taylor and perfect-power work. The following table locates several mathematically distinct routes beyond those developed in the main text.
\begin{center}\small
\begin{tabular}{@{}>{\raggedright\arraybackslash}p{.27\linewidth}>{\raggedright\arraybackslash}p{.20\linewidth}>{\raggedright\arraybackslash}p{.46\linewidth}@{}}\toprule
Earlier criterion & Family & Arithmetic content and location\\\midrule
Prime splitting & All three & One residue class, or primes of complete splitting, with the conductor or field allowed to grow. \code{SplitPrimes.lean}.\\[3pt]
Mertens and Liouville sums & RH, progressions for GRH, Euler coefficients for ERH & Integer power inequalities with growth witnesses. Perfect-power convolution and inversion convert related bounds. \code{MertensSyntax.lean}, \code{LiouvilleCriterion.lean}.\\[3pt]
Divisibility determinants & All three with the appropriate coefficients & First-column determinant expansion encodes summatory inverse coefficients. \code{DivisorMatrices.lean}, \code{MatrixIdentity.lean}.\\[3pt]
Farey discrepancies & Ordinary RH & A quadratic Farey energy is expressed using gcd sums and controlled with Mertens witnesses. \code{FareyRH.lean}, \code{EnergyInequality.lean}.\\[3pt]
Divisor sums & Ordinary RH & Weak and strict versions of Lagarias's $\sigma(n)\le H_n+e^{H_n}\log H_n$. \code{Lagarias.lean}.\\[3pt]
Witness transformations & All three & Arithmetic slack, prefix totals, and first finite failures have explicit derivations in object iPRA. \code{SlackForms.lean}, \code{ObjectForms.lean}, \code{CubeTests.lean}.\\\bottomrule
\end{tabular}
\end{center}
The Mertens formulation, for example, has the familiar RH meaning $\forall\epsilon>0\ |\sum_{n\le x}\mu(n)|=O_\epsilon(x^{1/2+\epsilon})$. The object formula uses integer powers $|M(N)|^{2(k+1)}\le C N^{k+2}$ with $\forall k\exists C\forall N$. The positive witness $C$ is part of this reading. In the Dedekind case the analogous summand comes from the reciprocal local Euler factors; it is not the ordinary M\"obius function on rational integers in every field. The companion report gives the field and residue-class bridges separately.

\paragraph{ProveIt sources.}
The submitted copies of Syntax, Deduction, PR and Semantics have Git blob hashes beginning \code{a3c8ec97}, \code{0969040f}, \code{32516315} and \code{c39a7004}, respectively. These agree with the corresponding files on the repository's main branch when checked. ProveIt's separate contour project specifies rational error, clearance and winding data, and describes its analytic interface as uncompiled. The prime-power approach here uses the ordinary iPRA presentation directly. See \href{https://github.com/VladimirReshetnikov/ProveIt/tree/ce0e3ceaf41a83e653857e5f60e0460733cc230c/Logic/PrimitiveRecursiveArithmetic/ConciseIPRA}{the pinned arithmetic framework} and the \href{https://github.com/VladimirReshetnikov/ProveIt/blob/main/Logic/PrimitiveRecursiveArithmetic/docs/proofs/analytic-number-theory/zeta_zeros_ipra_2026-11-15/lean/ObjectTargets.lean}{contour proof specification}.

\subsection*{The ambient analytic statements}
This source-level Lean listing uses Mathlib's continued Dirichlet $L$-function and its Dedekind zeta series. The listing was reviewed against the pinned \href{https://github.com/leanprover-community/mathlib4/blob/81a5d257c8e410db227a6665ed08f64fea08e997/Mathlib/NumberTheory/LSeries/DirichletContinuation.lean}{Dirichlet} and \href{https://github.com/leanprover-community/mathlib4/blob/81a5d257c8e410db227a6665ed08f64fea08e997/Mathlib/NumberTheory/NumberField/DedekindZeta.lean}{Dedekind} sources. Its analytic imports were unavailable as compiled dependencies for this build.
{\footnotesize\begin{verbatim}
def StripLine (f : Complex -> Complex) : Prop :=
  forall s : Complex, 0 < s.re -> s.re < 1 ->
    f s = 0 -> s.re = 1 / 2

def AmbientRH : Prop := RiemannHypothesis

def AmbientGRH : Prop :=
  forall (q : Nat) [NeZero q] (chi : DirichletCharacter Complex q),
    chi.IsPrimitive -> StripLine (DirichletCharacter.LFunction chi)

def Continuation (K : Type) [Field K] [NumberField K]
    (f : Complex -> Complex) : Prop :=
  And (Differentiable Complex f)
    (forall s : Complex, 1 < s.re ->
      f s = (s - 1) * NumberField.dedekindZeta K s)

def AmbientERH : Prop :=
  forall (K : Type) [Field K] [NumberField K],
    Exists (fun f => And (Continuation K f) (StripLine f))
\end{verbatim}}
Mathlib's \code{RiemannHypothesis} excludes negative even zeros and the designated value at $1$. The Dirichlet function in the display is already continued. The Dedekind series is used only on $\Re s>1$. Classical continuation provides the entire $f$, and uniqueness follows from the identity theorem on this half-plane. Multiplication by $s-1$ preserves strip zeros. Every number field has a small isomorphic copy, as required by the type universe above.

The compiled module \code{StripReflection.lean} proves the half-strip reduction. A functional equation propagates a zero to the reflected point, possibly in the dual character. The identity $\Re(1-s)=1-\Re s$ then yields \code{line_iff_right} and \code{line_iff_left}. These proofs apply to any family equipped with the stated reflection.


\paragraph{Further research questions.}
Can the one-sided Mellin proof be replaced by bounded finite inequalities, giving an object derivation from the sparse upper test back to the two-sided test? Can the Cayley translation bounds be combined with verified PR approximations to analytic continuation on a common domain for every field? For the residue-vector criterion, can a supplied prime-discrepancy bound be converted into explicit coefficients with a primitive recursive norm-error rate? 

\end{document}
